\section{Capítulo~\ref{sec:debugging}. Depuración de la máquina}

Qué oye el electrodo, la baja dimensión de la actividad, el funcionamiento de las interfaces con matrices rígidas, el aprendizaje conjunto del cerebro y el decodificador y los puntos visuales producidos por estimulación están medidos en experimentos revisados por pares; las estimaciones de los flujos de datos son aritmética del capítulo; sobre el dispositivo de Neuralink implantado en personas, hasta donde pudimos comprobar, los datos los ha publicado sobre todo la propia empresa.

{\footnotesize
\setlength{\tabcolsep}{3pt}\renewcommand{\arraystretch}{1.15}
\begin{longtable}{>{\raggedright}p{0.5cm}>{\raggedright}p{4.0cm}>{\raggedright}p{5.6cm}>{\raggedright}p{2.3cm}>{\raggedright\arraybackslash}p{1.7cm}}
\toprule
N.º & Afirmación & Experimento y qué se midió & Fuente & Estado \\
\midrule\endfirsthead
\toprule
N.º & Afirmación & Experimento y qué se midió & Fuente & Estado \\
\midrule\endhead
1 & El electrodo oye las espigas de las neuronas en un radio del orden de un centenar de micrómetros, normalmente de una a unas pocas; se distinguen por su forma & Rata, área CA1, registro intracelular y extracelular simultáneo: la forma de la espiga extracelular reproduce la anchura y la amplitud de la intracelular. Estimación: un tetrodo podría distinguir $60$--$100$ neuronas, pero en la práctica se aíslan unas seis. & \cite{Henze2000extra} & medido \\
2 & El cerebro tiene $8.6\cdot10^{10}$ neuronas; la sonda oye más o menos una cienmillonésima & Recuento de núcleos en una suspensión de tejido disuelto: $86\cdot10^9$ neuronas en el cerebro humano. La fracción es aritmética del capítulo. & \cite{Azevedo2009} & medido \\
3 & La actividad de cientos de neuronas de la corteza motora cabe en una o dos decenas de variables comunes & Revisión de registros en la corteza motora de monos: la actividad de la población se describe con pocas variables latentes comunes a muchas neuronas. & \cite{Gallego2017} & medido \\
4 & El ruido de las vecinas está correlacionado, y un centenar de neuronas lleva casi lo mismo que mil (proposición~\ref{prop:pooling}) & Corteza visual de mono: coeficiente de correlación de ruido entre neuronas vecinas de alrededor de $0.1$, y la ganancia por promedio se satura con un centenar de neuronas. Teoría de las correlaciones que limitan la información. & \cite{Zohary1994, MorenoBote2014} & medido \\
5 & Flujo bruto de $2\cdot10^8$~bit/s frente a $8\cdot10^4$~bit/s en espigas~\eqref{eq:probedata} & Aritmética del capítulo basada en la capacidad del flujo de espigas~\eqref{eq:spikeentropy}. Los parámetros de digitalización son reales: el chip de Neuralink trabaja a $19.3$~kHz y $10$~bits por canal. & deducción del capítulo; \cite{Rieke1997, Musk2019} & teoría \\
6 & Primer sistema de Neuralink: los chips amplifican y digitalizan, el flujo bruto va por cable y las espigas las detecta un programa externo; la detección en el chip, la carcasa cerrada, la radio y la alimentación inalámbrica corresponden al dispositivo para personas, según la empresa & Artículo de la empresa: todo el flujo bruto por un cable USB-C; las espigas las detecta en tiempo real un programa fuera del chip; la compresión a bordo, la alimentación y la transmisión inalámbricas figuran como plan para los dispositivos clínicos. Para el dispositivo implantado en personas, solo comunicados de la empresa. Que sea necesario detectar las espigas en el chip es una conclusión del capítulo a partir de~\eqref{eq:probedata}. & \cite{Musk2019}; informes de Neuralink, sin artículo revisado por pares & medido (artículo de la empresa); en personas, informe de la empresa \\
7 & Hasta $3072$ electrodos en $96$ hilos, $32$ por hilo; los hilos los inserta un robot, esquivando los vasos & Artículo de la empresa: $3072$ electrodos en $96$ hilos de $32$; el robot inserta $6$ hilos ($192$ electrodos) por minuto; hasta el $70\%$ de los electrodos oyen espigas en ratas con la matriz implantada a largo plazo. & \cite{Musk2019} & medido (artículo de la empresa) \\
8 & En personas, unos mil canales en $64$ hilos; parte de los hilos se desplazó; cursor del orden de $10$~bit/s & Solo comunicados de la empresa. Una búsqueda en PubMed no encontró ningún artículo revisado por pares con resultados en personas; esto concuerda con la salvedad del texto del capítulo. & informes de Neuralink; sin artículo revisado por pares & medido (informe de la empresa) \\
9 & Una interfaz con una matriz de un centenar de electrodos controla un cursor & Persona con parálisis, $96$ electrodos en la corteza motora primaria: la intención de mover el brazo modifica las espigas tres años después de la lesión; el decodificador da un “cursor neuronal” (correo, televisor), control de una prótesis de mano y de un brazo robótico. & \cite{Hochberg2006} & medido \\
10 & Escritura con la “mano mental”: unos $90$ caracteres por minuto, menos de $8$~bit/s & Persona con parálisis de la mano, intentos de escribir a mano, decodificador recurrente: $90$ caracteres por minuto con una precisión del $94.1\%$ en tiempo real, más del $99\%$ tras la autocorrección. La estimación en bits es aritmética del capítulo. & \cite{Willett2021} & medido \\
11 & Los intentos de hablar se convierten en texto a unas $60$ palabras por minuto & Persona que perdió el habla inteligible, matrices en la corteza: $62$ palabras por minuto; $9.1\%$ de errores de palabra con un vocabulario de $50$ palabras y $23.8\%$ con uno de $125\,000$. & \cite{Willett2023} & medido \\
12 & La interfaz extrae una pequeña fracción de la capacidad: una neurona, decenas de bits por segundo; un centenar, miles (proposición~\ref{prop:probe}) & El límite superior~\eqref{eq:spikeentropy} es teoría; la comparación con las filas~8 y~10 es una conclusión del capítulo. Que el cuello de botella sea la baja dimensión y el desconocimiento del código, y no el cable, no se ha comprobado con un experimento directo. & \cite{Rieke1997} & teoría \\
13 & El cerebro y el decodificador aprenden el uno del otro & Monos controlan un cursor; el decodificador se ajusta en lazo cerrado mientras las neuronas se reaprenden a la vez: la precisión aumenta, la habilidad se conserva y le afectan menos los movimientos del propio brazo. El consejo de separar las velocidades de aprendizaje es una regla de ingeniería, sin experimento propio en el capítulo. & \cite{Orsborn2014coad} & medido \\
14 & La corriente a través de un electrodo excita las neuronas y los cables circundantes sin distinguir tipo & Corteza de ratón y de gato, registro de calcio con dos fotones: incluso una corriente débil excita neuronas escasas, dispersas hasta a milímetros de distancia, mediante la excitación directa de axones en un volumen de decenas de micrómetros de diámetro. Es decir, la excitación no es selectiva, pero tampoco continua. & \cite{Histed2009stim} & medido \\
15 & La estimulación de la corteza visual enciende puntos; hasta $16$ puntos a la vez permiten reconocer algunas letras y los bordes de objetos & Persona ciega, $96$ electrodos en la corteza visual durante $6$ meses: umbral de un electrodo, $66.8\pm36.5$~\textmu A; la estimulación simultánea de grupos (hasta $16$ electrodos, el límite del estimulador) reduce el umbral y da patrones distinguibles; se reconocen algunas letras y bordes de objetos. & \cite{Fernandez2021} & medido \\
16 & La segunda línea de Neuralink es un dispositivo que introduce señales en la corteza visual & Solo comunicados de la empresa sobre su desarrollo; no se encontraron datos revisados por pares sobre su funcionamiento. & informes de Neuralink & hipótesis \\
17 & Las concentraciones de los neurotransmisores de difusión pueden medirse con un sensor químico en el tejido & Rodajas de cerebro de rata, voltamperometría cíclica rápida con fibra de carbono: se midieron la liberación y la recaptación de serotonina; la sustancia sale fuera de la sinapsis. & \cite{Bunin1998} & medido \\
\bottomrule
\end{longtable}}
